\subsection{Finite-fibre transfer strictness}

\begin{proposition}[Finite-fibre certificates and transfer obstruction]
\label{prop:L2-generic-vs-structural}
Finite-fibre certificates are fixed-resolution finite-atomic data.  A
finite-state transfer formula in the resolution variable is extra structure.
\end{proposition}

\begin{proof}
This is \Cref{prop:L2-novelty-finite-state-separation}.
\end{proof}

\begin{definition}[Computable finite-fibre cover]
\label{def:computable-cover}
A computable finite-fibre cover of the Zeckendorf resolution data is auxiliary
finite certificate data over the fixed normal-form set \(X_m\): over each normal
form one attaches a computably specified nonempty finite set, without changing
\(X_m\), the original maps \(\Fold_m\), or any readout on normal forms.
\end{definition}

\begin{lemma}[Covers preserve readouts]
\label{lem:cover-preserves-protocol}
Computable finite-fibre covers preserve the Zeckendorf normal forms and the
compiled \(L1\) readouts.
\end{lemma}

\begin{proof}
The cover changes only the number of certificate points over a fixed normal
form.  Projection to the normal form is unchanged, so every operation and
readout defined on normal forms is unchanged.
\end{proof}

\begin{theorem}[Certificate--transfer dichotomy for \texorpdfstring{\(L2\)}{L2}]
\label{thm:L2-finite-state-novelty-spine}
Fix the Zeckendorf normal-form system \((X_m,\Fold_m)\).  Finite-fibre
certificate universality at level \(L2\) is fixed-resolution finite-atomic data;
a fixed finite-state transfer law in the resolution parameter is additional
cross-resolution structure.  More precisely, in the computable finite-fibre
cover category of \Cref{def:computable-cover}:
\begin{enumerate}
\item a cover is exactly a computable positive integer fibre-size function
  \(d_m:X_m\to\mathbf N_{\ge1}\), and projection to \(X_m\) preserves all
  Zeckendorf normal forms and compiled \(L1\) readouts;
\item if the second covered collision moments
  \[
    M_m(d)=\sum_{x\in X_m} d_m(x)^2
  \]
  are governed by a fixed finite-state transfer formula, then
  \(M_m(d)\le C\lambda^m\) for some constants \(C,\lambda>0\);
\item for every finite window \(N\), every rational \(A>0\), and every rational
  \(\rho>1\), there is a computable cover, equal to the singleton cover for
  \(m<N\), whose moments satisfy \(M_m(d)>A\rho^m\) for all \(m\ge N\) and are
  not governed by any fixed finite-state transfer formula.
\end{enumerate}
Thus the \(L2\) novelty proved here is finite-fibre certificate universality;
finite-state transfer is a rigid strengthening that can be imposed for the
intrinsic fold moments of \Cref{thm:intrinsic-fold-moment-transfer} but is not
forced by finite-fibre certificates or by any finite initial certificate window.
This is the certificate/transfer boundary used by
\Cref{thm:finite-state-strictness,thm:L2-cover-relative-closure} in the
\(L2\) clause of \Cref{thm:main-hierarchy}.
\end{theorem}

\begin{proof}
For a fixed resolution, \Cref{prop:classical-finite-atomic-skeleton} identifies
finite-fibre certificates with finite atomic certificate data.  In the cover
category this datum is exactly a positive integer fibre-size function over each
normal form: one direction counts the added finite certificate fibre over
\(x\in X_m\), and the other attaches \(\{1,\ldots,d_m(x)\}\) over \(x\).  The
constructions are computable and inverse up to relabelling of certificate
points.  The preservation of normal forms and \(L1\) readouts is
\Cref{lem:cover-preserves-protocol}.  This proves the first assertion.

For the second assertion, a fixed finite-state transfer formula has only
finitely many states and fixed rational transition and output data.  By
\Cref{prop:finite-state-coefficient-budget,prop:fixed-transfer-interface-budget},
every sequence produced by such a fixed transfer formula is bounded by
\(C\lambda^m\) in the resolution parameter for suitable constants \(C\) and
\(\lambda\).

It remains to prove the quantitative non-forcing assertion.  Choose a computable
distinguished normal form \(x_m\in X_m\), for instance the lexicographically first
normal form.  Put \(d_m\equiv1\) for \(m<N\).  For \(m\ge N\), put
\(d_m(x)=1\) for \(x\ne x_m\) and
\[
  d_m(x_m)=1+m!+\left\lceil \sqrt{A\rho^m}\right\rceil .
\]
This is a computable finite-fibre cover of the unchanged normal-form system.
For \(m\ge N\), its second covered moment is
\[
  M_m(d)=\bigl(1+m!+\lceil\sqrt{A\rho^m}\rceil\bigr)^2+|X_m|-1>A\rho^m.
\]
Moreover \(M_m(d)\ge (m!)^2\) for all large \(m\).  If the sequence had a fixed
finite-state transfer formula, the previous paragraph would give
\(M_m(d)\le C\lambda^m\) for some \(C,\lambda\), contradicting the elementary
fact that \((m!)^2/(C\lambda^m)\to\infty\).  Hence no such transfer formula
exists.  Since the construction is the singleton cover on the prescribed window
and every cover projects to the same unmodified \((X_m,\Fold_m)\), finite-state
transfer is not determined by finite-fibre certificates, by finite prefix data,
or by the intrinsic fold map without an additional comparison theorem.  The
positive intrinsic transfer theorem is therefore a separate strengthening, not
the definition of \(L2\).
\end{proof}

\begin{theorem}[Computable covers are computable fibre-size data]
\label{thm:finite-fiber-cover-realization}
Specifying a computable finite-fibre cover is equivalent to specifying a
computable positive integer fibre size over each normal form at each resolution.
\end{theorem}

\begin{proof}
Given a cover, count the finite certificate fibres effectively.  Given
computable positive sizes, attach the auxiliary certificate fibre
\(\{1,\ldots,d\}\) over the corresponding normal form.  These
constructions are inverse up to relabelling of certificate points.
\end{proof}

\begin{theorem}[Extremal cover spine used for \texorpdfstring{\(L2\)}{L2} strictness]
\label{thm:L2-cover-skeleton-strictness}
Fix the Zeckendorf normal-form system \((X_m,\Fold_m)\).  In the class of
computable finite-fibre covers of \Cref{def:computable-cover}, the
finite-state-transfer obstruction proved in this article is witnessed, and at
fixed certificate excess is maximized, by the computable fibre-size function over
the fixed sets \(X_m\).  More precisely, the cover-relative strictness theorem
proved here follows from the following three-part spine:
\begin{enumerate}
\item finite-state transfer sequences are exponentially bounded in the
  resolution parameter by the fixed-transfer budget of
  \Cref{prop:fixed-transfer-interface-budget}; and
\item among covers with a prescribed total excess \(E_m\) above singleton fibres
  at resolution \(m\), the second covered collision moment is at most
  \((E_m+1)^2+|X_m|-1\), with equality precisely for a one-spike cover; and
\item a computable cover may realize this extremal one-spike fibre over a fixed
  normal form at resolution \(m\) for any prescribed computable excess \(E_m\).
\end{enumerate}
Moreover, the same one-spike skeleton can be chosen to agree with any prescribed
finite certificate window before it violates the transfer budget.
No assertion about the intrinsic fibre multiplicities of the original maps
\(\Fold_m\) is used or implied by this skeleton.
\end{theorem}

\begin{proof}
The first part is \Cref{prop:fixed-transfer-interface-budget}: by
\Cref{def:fixed-finite-state-transfer,prop:finite-state-coefficient-budget}, any
fixed finite-state transfer formula gives a sequence bounded by
\(C\lambda^m\) for some constants \(C,\lambda\).  For the extremal assertion,
write the cover sizes at resolution \(m\) as \(1+e_x\), where
\(e_x\ge0\) and \(\sum_x e_x=E_m\).  Then
\[
  \sum_x(1+e_x)^2=|X_m|+2E_m+\sum_xe_x^2
  \le |X_m|+2E_m+E_m^2=(E_m+1)^2+|X_m|-1,
\]
and equality holds exactly when all nonzero excess is concentrated at one normal
form.  For the realizing cover, use \Cref{thm:finite-fiber-cover-realization}:
choose one normal form in each
nonempty \(X_m\), assign its cover fibre size \(E_m+1\), and assign size \(1\) to
all other cover fibres.  This is a computable finite-fibre cover, and by
\Cref{lem:cover-preserves-protocol} the projection to \(X_m\), the normal forms,
and all \(L1\) readouts are unchanged.  Its second collision moment is therefore
the extremal value \((E_m+1)^2+|X_m|-1\); taking a superexponential computable
sequence such as \(E_m=m!\) violates the exponential budget and gives
\Cref{thm:finite-state-strictness}.  If an initial certificate window is compared
with the singleton cover, take fibre size \(1\), equivalently excess \(E_m=0\),
on that window and make the same superexponential choice afterwards; this is
exactly the finite-window obstruction of
\Cref{thm:L2-finite-window-transfer-obstruction}.  Since the construction
changes only cover certificate multiplicities above already fixed normal forms,
it makes no claim about the unmodified fibres of \(\Fold_m\).
\end{proof}

\begin{theorem}[Quantitative finite-state budget obstruction]
\label{thm:L2-quantitative-budget-obstruction}
Let \(A>0\) and \(\rho>1\) be rational numbers, and let \(N\ge1\).  There is a
computable finite-fibre cover of the Zeckendorf normal-form system with the
following properties.
\begin{enumerate}
\item For every \(m<N\) it is the singleton cover, so its finite-fibre
  certificate multiset and all \(L1\) readouts agree with the unmodified
  normal-form presentation at those resolutions.
\item For every \(m\ge N\), its second covered collision moment is larger than
  \(A\rho^m\).
\item The cover changes only one auxiliary certificate fibre over the fixed set
  \(X_m\) at each resolution.  Hence the budget violation is purely
  cover-relative and gives no assertion about the intrinsic multiplicities of
  \(\Fold_m\).
\end{enumerate}
Consequently no finite initial certificate window, and no effective rational
exponential finite-state transfer budget, can characterize \(L2\) finite-fibre
certificates inside the computable cover category.
\end{theorem}

\begin{proof}
For \(m<N\), put singleton fibres over every element of \(X_m\).  For
\(m\ge N\), choose one normal form \(x_m^\ast\in X_m\) and put
\[
  E_m=\left\lceil \sqrt{A\rho^m}\right\rceil,
\]
then attach a fibre of size \(E_m+1\) over \(x_m^\ast\) and singleton fibres over
all other normal forms.  The function \(m\mapsto E_m\) is computable from the
fixed rational constants as an integer-valued sequence; equivalently, compute
\(E_m\) as the least integer \(E\) with \(E^2\ge A\rho^m\).  By
\Cref{thm:finite-fiber-cover-realization} this fibre-size function defines a
computable finite-fibre cover, and by \Cref{lem:cover-preserves-protocol} it
preserves all normal forms and all \(L1\) readouts.

For \(m<N\), the construction is singleton by definition.  For \(m\ge N\), the
second covered moment is
\[
  (E_m+1)^2+|X_m|-1.
\]
Since \(|X_m|\ge1\) and \(E_m^2\ge A\rho^m\), this is strictly larger than
\(A\rho^m\).  Only the single auxiliary fibre over \(x_m^\ast\) is enlarged;
the original sets \(X_m\) and maps \(\Fold_m\) are not changed.  Thus the
construction defeats the specified exponential transfer budget while remaining
entirely in the cover category.
\end{proof}

\begin{theorem}[Main-article boundary for finite-fibre refinements]
\label{thm:L2-main-article-boundary}
Only computable finite-fibre cover data that leave normal forms fixed are used
in this article's \(L2\) strictness theorem.  The unmodified Zeckendorf
\(\Fold_m\) fibre multiplicities are covered instead by the intrinsic transfer
and height package of
\Cref{thm:intrinsic-fold-moment-transfer,thm:intrinsic-second-moment-recurrence,cor:polynomial-intrinsic-fiber-statistics,thm:intrinsic-fold-fibre-height,cor:linfty-fold-entropy,thm:zero-temperature-fold-pressure}.
No full intrinsic fibre-histogram law or large-deviation theorem is used in the
cover strictness chain.
\end{theorem}

\begin{proof}
All results here use \Cref{def:computable-cover}, whose projection to
\(X_m\) is fixed.  The cover-size function is independent finite certificate data
over the already chosen normal forms, while the intrinsic multiplicities are the
fibre sizes of the unmodified maps \(\Fold_m\).  By
\Cref{thm:L2-cover-intrinsic-nonimplication}, a transfer statement or obstruction
for the cover moment has no consequence for the intrinsic fold moment without an
additional comparison theorem.  The intrinsic fold moment and height endpoint are
handled directly by
\Cref{thm:intrinsic-fold-moment-transfer,thm:intrinsic-fold-fibre-height,cor:linfty-fold-entropy,thm:zero-temperature-fold-pressure}; the cover strictness proof supplies no
full intrinsic histogram or large-deviation classification.
\end{proof}

\begin{theorem}[Computable-cover-only boundary of finite-state obstruction]
\label{prop:cover-only-strictness-boundary}
\label{thm:one-spike-finite-state-boundary}
The one-spike construction is sharp for the computable finite-fibre cover model:
it preserves all normal-form readouts and changes only one added finite
certificate fibre at each resolution, yet it can realize any computable
second-moment spike of the form \((a_m+1)^2+F_{m+2}-1\).  The theorem is silent
about intrinsic \(\Fold_m\) fibre multiplicities.
\end{theorem}

\begin{proof}
This is the construction in \Cref{thm:L2-finite-fiber-no-transfer-content}.
Only one fibre is enlarged, and \Cref{lem:cover-preserves-protocol} proves that
readouts are preserved.
\end{proof}

\begin{theorem}[Intrinsic height is separate from cover spikes]
\label{thm:intrinsic-fold-asymptotics-split}
The factorial one-spike obstruction does not assert a full intrinsic histogram or
large-deviation law for the original Zeckendorf fibres.  It is a statement about
computable finite-fibre covers, so no such classification theorem for
\(\Fold_m\) fibre multiplicities follows from
\Cref{thm:finite-state-strictness}.  The positive intrinsic fixed-\(q\) transfer,
recurrence, height, and zero-temperature endpoint statements used in the article
are exactly those proved in
\Cref{thm:intrinsic-fold-moment-transfer,thm:intrinsic-second-moment-recurrence,cor:polynomial-intrinsic-fiber-statistics,thm:intrinsic-fold-fibre-height,cor:linfty-fold-entropy,thm:zero-temperature-fold-pressure}.
\end{theorem}

\begin{proof}
The proof of \Cref{thm:L2-finite-fiber-no-transfer-content} deliberately alters
cover multiplicities.  It leaves the original fold map untouched.  The formal
nonimplication is \Cref{thm:L2-cover-intrinsic-nonimplication}.
\end{proof}


\begin{theorem}[Cover-relative closure of the \texorpdfstring{\(L2\)}{L2} strictness clause]
\label{thm:L2-cover-relative-closure}
The \(L2\) strictness clause used by \Cref{thm:main-hierarchy} is fully proved
inside the computable finite-fibre cover category of \Cref{def:computable-cover}.
It neither requires nor implies any intrinsic nonrecurrence, finite-state-failure,
or asymptotic growth classification for the multiplicity sequence of the original
maps \(\Fold_m\).  Positive intrinsic finite-state transfer for those original
maps is supplied separately by \Cref{thm:intrinsic-fold-moment-transfer}.
\end{theorem}

\begin{proof}
The positive finite-fibre part of \(L2\) is fixed-resolution certificate data:
\Cref{prop:classical-finite-atomic-skeleton,prop:zeckendorf-cert-univ} give the
finite-atomic certificates from finite fibres at each fixed resolution.  A
finite-state transfer law is additional cross-resolution
structure.  By
\Cref{def:fixed-finite-state-transfer,prop:fixed-transfer-interface-budget},
every fixed finite-state transfer sequence has an exponential bound in the
resolution parameter.

Now work in the cover category.  By \Cref{thm:finite-fiber-cover-realization}, a
computable cover is exactly a computable positive integer fibre-size function
over the unchanged normal-form sets \(X_m\), and by
\Cref{lem:cover-preserves-protocol} such a cover preserves the Zeckendorf normal
forms and all \(L1\) readouts.  The one-spike skeleton of
\Cref{thm:L2-cover-skeleton-strictness} gives the extremal obstruction at each
fixed certificate excess, and the quantitative budget form
\Cref{thm:L2-quantitative-budget-obstruction} says more: after any prescribed
finite window, and against any effective rational exponential bound \(A\rho^m\), a
computable one-spike cover preserves the protocol while forcing the second
covered moment above that bound.  In particular, taking a superexponential
computable spike such as \(m!\) gives a computable cover whose second collision
moment violates every exponential bound, and therefore cannot come from a fixed
finite-state transfer formula in the sense of \Cref{def:fixed-finite-state-transfer}.
This proves the cover-relative strictness theorem.

The proof has not used the intrinsic fibre sizes of the original maps
\(\Fold_m\).  Indeed \Cref{thm:L2-cover-intrinsic-nonimplication} records the
formal nonimplication: altering auxiliary cover multiplicities over fixed normal
forms gives no comparison theorem for unmodified fold fibres.  Hence the main
article's \(L2\) strictness is exactly cover-relative.  The positive intrinsic
transfer package is the separate theorem package in \Cref{sec:L2}, not a
consequence of the one-spike cover obstruction.
\end{proof}

\begin{theorem}[Cover-relative finite-fibre certificates do not force finite-state transfer]
\label{thm:finite-state-strictness}
\label{thm:cover-relative-finite-fibre-certificates-do-not-force-finite-state-transfer}
There is a computable finite-fibre cover of the Zeckendorf \(L1\) protocol whose
finite-fibre certificates exist at every resolution but whose second collision
moment over the fixed normal-form sets \(X_m\) has no fixed finite-state transfer
formula in the sense of \Cref{def:fixed-finite-state-transfer}.  This is a
cover-relative strictness theorem in the computable
finite-fibre cover category; it is not an intrinsic finite-state failure theorem
for the original \(\Fold_m\) fibre multiplicities and does not contradict the
intrinsic height law of \Cref{thm:intrinsic-fold-fibre-height}.  More sharply,
the obstruction is quantitative: after any prescribed finite certificate window
and against any effective rational exponential transfer budget, a one-spike computable
cover can preserve all normal forms and all \(L1\) readouts while forcing the
second covered moment beyond that budget.  Over the same unmodified Zeckendorf
data \((X_m,\Fold_m)\) there are therefore computable covers whose second
collision moments have opposite finite-state-transfer behaviour, as isolated in
\Cref{prop:cover-category-invariance-finite-state-obstruction}.  Hence the
obstruction is not a property determined by the intrinsic fold map alone.
\end{theorem}

\begin{proof}
Apply the cover skeleton \Cref{thm:L2-cover-skeleton-strictness} with
\(a_m=m!\), equivalently apply \Cref{thm:L2-finite-fiber-no-transfer-content} to
that one-spike sequence.  The cover is computable and preserves readouts by
\Cref{lem:cover-preserves-protocol}.  Its second moment grows faster than every
exponential, whereas every fixed finite-state transfer sequence is exponentially
bounded by \Cref{prop:fixed-transfer-interface-budget}.  The obstruction is
therefore cover-relative: it concerns the added finite certificate fibres over
the unchanged normal-form sets, not intrinsic multiplicities of the original
\(\Fold_m\) fibres.

The sharper two-cover assertion is
\Cref{prop:cover-category-invariance-finite-state-obstruction,thm:L2-two-cover-intrinsic-boundary}.
Its singleton cover has moment \(F_{m+2}\), which has a fixed two-state Fibonacci
transfer formula, while its factorial one-spike cover has moment
\((m!+1)^2+F_{m+2}-1\), which violates the exponential budget just cited.  Since
both covers project to the same unmodified pair \((X_m,\Fold_m)\), finite-state
transfer or failure for these cover moments cannot be read off from the intrinsic
fold map without an additional bridge identifying intrinsic fold multiplicities
with the chosen cover multiplicities.  Such a bridge is outside the theorem chain
by \Cref{thm:L2-cover-intrinsic-nonimplication}.
\end{proof}

\begin{proposition}[Finite prefixes do not certify transfer]
\label{prop:finite-prefix-no-transfer-certification}
In the computable finite-fibre cover category, no bounded initial window of
finite-fibre certificate data can certify finite-state transfer for all later
resolutions.
\end{proposition}

\begin{proof}
Given any bound \(M\), choose a one-spike cover with \(a_m=1\) for
\(m\le M\) and \(a_m=m!\) for \(m>M\).  It agrees with the unspiked data on the
initial window but violates the finite-state exponential budget later.
\end{proof}

\begin{theorem}[Budget spine used for \texorpdfstring{\(L2\)}{L2} strictness]
\label{thm:L2-budget-spine}
The cover-relative \(L2\) strictness obstruction proved in this article is a
growth-budget obstruction: finite-state transfer sequences obey the exponential
bound of \Cref{prop:fixed-transfer-interface-budget}, while computable
finite-fibre covers can violate it, and the superexponential cover obstructions
used here have the one-spike skeleton of
\Cref{thm:L2-cover-skeleton-strictness}.
\end{theorem}

\begin{proof}
The budget is \Cref{prop:fixed-transfer-interface-budget}, which applies the
definition \Cref{def:fixed-finite-state-transfer} through
\Cref{prop:finite-state-coefficient-budget}; the violating cover is
\Cref{thm:finite-state-strictness}.  The one-spike bad-example skeleton used for
the cover-relative strictness theorem is \Cref{thm:L2-cover-skeleton-strictness}.
\end{proof}

\begin{corollary}[Finite-state growth budget]
\label{cor:finite-state-growth-budget}
Every fixed finite-state transfer formula gives at most exponential growth in
the resolution parameter.
\end{corollary}

\begin{proof}
This is the exponential-bound clause of
\Cref{prop:fixed-transfer-interface-budget}.
\end{proof}



